\documentclass{zhvt-classic}

% 使用本地源流明体（GenRyuMin2 TW）按文件路径加载，不依赖系统字体名
% 版框（内框／界栏／书口侧栏／鱼尾）改用与版心印章同一印泥棕色：
%   印章图片 yshs.png 实色 = RGB(135,68,52) = HTML 874434
\usepackage{xcolor}
\usepackage{tikz}
% sealbrown 双注册：xcolor 供 TikZ/pgf，L3 供 cls 版框
\definecolor{sealbrown}{HTML}{874434}
\ExplSyntaxOn
\color_set:nnn { sealbrown } { HTML } { 874434 }
\ExplSyntaxOff

\zhvtset{
  grid color       = sealbrown,
  main font        = GenRyuMin2TW-SB,
  main font path   = C:/Users/Administrator/texmf/fonts/opentype/public/GenRyuMin2TW/,
jiazhu font      = GenRyuMin2TW-M,
  jiazhu font path = C:/Users/Administrator/texmf/fonts/opentype/public/GenRyuMin2TW/,
  % 页码框容纳 4 字：正常页「九十七」+前置空 1 格=4 行恰好；
  % 100+ 页「一百零一」共 5 行超出，末字「一/二/三」被框裁掉（2026-10-04 实测），故放宽到 5
  shukou page number chars = 5,
  % 书口书名框高从默认 6 字减到 4 字，把省出的 2 字高拨给 chapter 框，
  % 使 6 字书名+卷名放进 9 字高的 chapter 框后不再压线（2026-10-05）
  shukou title chars = 4,
}

\title{史记}[淮南衡山列传]
\author{汉·司马迁}
\maker{闲叟漫记·重製}


% ===== 索隐述赞样式（2026-11-08 四次定稿 · 愚大人亲手改定参数）=====
% ⚠ 竖排前提：自定义字体家族必须带 RawFeature={vertical:+vert}；
%   simhei.ttf 加 +vhal 会整体加载失败静默回退，只写 +vert；simfang 同理。
\newCJKfontfamily{\shuzanfont}[Path=C:/Windows/Fonts/,RawFeature={vertical:+vert}]{simhei.ttf}
\newCJKfontfamily{\shuzanbody}[Path=C:/Windows/Fonts/,RawFeature={vertical:+vert}]{simfang.ttf}
% 题头：152x30pt 圆角 6pt 棕底白字；baseline (0,2) 使整框落入界行格子不压版框线；
% 白字节点 x 不补偿、y +8.2pt（像素级墨迹实测 dx/dy≈0，L/R 留白约 20pt）。
\newcommand{\shuzanhead}[1]{%
  \begin{tikzpicture}[baseline={(0,2)}]
    \fill[sealbrown, rounded corners=6pt]
      (-76pt,-15pt) rectangle (76pt,15pt);
    \node[text=white, inner sep=0pt, anchor=center] at (0pt,8.2pt)
      {\shuzanfont\fontsize{23}{23}\selectfont #1};
  \end{tikzpicture}%
}
% ★ 修正书口卷名位置：chapter 框起始位置改回 cls 默认值，避免全局影响。
%   只在每个卷首页通过 \juan 命令里的空命令微调该页位置（安全，不会跨页污染）。
\ExplSyntaxOn
\AtBeginDocument{
  \dim_set:Nn \l__zhvt_shukou_chapter_title_start_dim
    { \l__zhvt_shukou_title_height_dim + \l__zhvt_font_size_dim * 2 }
  \dim_set:Nn \l__zhvt_shukou_page_number_start_dim
    { \l__zhvt_shukou_chapter_title_start_dim + \l__zhvt_shukou_chapter_title_height_dim }
}
\ExplSyntaxOff

\begin{document}

% ★★ 版心（书口）排布覆写：把书名放到鱼尾之下（2026-10-04）
%
%   cls 默认版心自上而下：书名（title 静态框，在鱼尾之上）→ 鱼尾 → 篇名（chapter 动态框）→ 页码
%   本模板改为：          鱼尾（位置不动）→ 书名 → 篇名 → 页码分割线 → 页码
%
%   手法：把 cls 放在 title 静态框里的书名，改放进 chapter 动态框（鱼尾之下）；
%         title 静态框留空。书名取 \l__zhvt_title_str，即 \title{...} 的第一参数，
%         所以换书时只改 \title，不必在此重复写书名。
%
%   ⚠ 三条必须遵守的写法（2026-10-04 实测）：
%     ① 汉字必须用 \hbox{...} 包住。flowfram 会把框内容装进一个宽 14.4mm（= 一个 line_sep
%        = 书口栏宽）的纵向 minipage 并按该宽度折行，而整页靠 PDF /Rotate 90 显示、
%        框内文字本身不变向 —— 不包 \hbox 时「围炉夜话」4 字会被逐字折成 4 行，
%        在版面上就成了横排 4 字并横向溢出版框。
%     ② 竖排下**一个 \hbox 就是一整列**。书名与篇名要同处一列，必须放进同一个 \hbox；
%        分成两个 \hbox 时，后一个会被当成新的一列、甩到隔壁栏去
%        （实测「序」跑到书口栏外侧 14.4mm 的正文列上）。中间的 \quad 在框内是横向空白，
%        显示出来正好是书名与篇名之间的纵向间隔。
%     ③ ★ 内容盒要写成 **与框等宽的定宽盒** `\hbox to \l__zhvt_shukou_chapter_title_height_dim{…\hfil}`。
%        实测：chapter 框是把内容盒沿纵向**居中**摆放的，内容宽度一变，书名就跟着漂：
%        自然宽度的 \hbox（书名 38.4mm）落在 x=157.2mm，加篇名后（57.6mm）落在 x=152.4mm，
%        目录页（无篇名）与正文页的书名就差半格。做成等宽盒后，无论有没有篇名，
%        内容固定从框顶起排，实测两页均落在 x=142.8mm，不再漂。
\ExplSyntaxOn
% 把书口 chapter 框（书名+篇名）起始位置下移 1 字高，避免 4 字卷名顶到鱼尾
% ⚠ cls 的 \__zhvt_calculate_dimensions: 在 \AtBeginDocument 里执行，会重算该 dim；
%   必须把覆写挂到 begin document 之后，否则会被二次覆盖（2026-10-05 实测）。
% (旧方案已移到导言区，直接修改 chapter 框起始位置)
% 卷标题命令 #1=书口篇名(鱼尾) #2=正文卷名（含卷序，卷名内空格用 ~）
%
% ★★ 目录页两项修订（2026-10-05，愚大人指出）★★
%   ① 页码起点统一对齐：
%      原状 —— 目录项写成 \par\hspace{4em}卷名\hspace{4em}页码，页码起点 = 卷名实宽 + 4em。
%      卷名实宽随字数变：「卷一~梁惠王上」6 字 +~、「卷七~离娄上」5 字 +~，
%      差 1 个字位 → 卷七~卷十 4 条的页码整体比其余 10 条高 27.21pt（一个字），看着上下不齐。
%      改法 —— 卷名补全角空格（宽 = 1 字 = 27.215pt）补足到 **7 字符位**再写进目录，
%      14 条页码起点遂全部落在同一条竖线上（起点对齐；末字随页码字数自然参差）。
%      ⚠ 只对目录里的副本 \l_tmpc_tl 补，正文卷名仍用 \l_tmpa_tl 原样 —— 正文像素不受影响。
%   ② 卷名整体上移半字位：前导 \hspace 由 4em 改 3.5em。
%      实测原状目录块在版框内偏下（天头留白 139.6pt / 地脚留白 110.7pt），
%      上移 0.5em = 13.6pt 后上下基本均分，视觉居中。
\NewDocumentCommand { \juan } { mm } {
	\clearpage
	\tl_set:Nn \l__zhvt_chapter_title_tl { #1 }
	\tl_set:Nn \l_tmpa_tl { #2 }
	\tl_set_eq:NN \l_tmpc_tl \l_tmpa_tl
	\int_step_inline:nn { 7 - \str_count:n { #2 } }
	  { \tl_put_right:Nn \l_tmpc_tl { 　 } }
	\tl_set:Nx \l_tmpb_tl { \zhnum { page } }
	\tl_put_left:Nn \l_tmpb_tl { \hspace { 4em } }
	\tl_put_left:Nn \l_tmpb_tl { \l_tmpc_tl }
	\tl_put_left:Nn \l_tmpb_tl { \par \hspace { 3.5em } }
	\addtocontents { toc } { \l_tmpb_tl }
	\skip_horizontal:n { \l__zhvt_font_size_dim * 4 }
	\l_tmpa_tl
	\par
	\thispagestyle { empty }
}
\cs_set_protected:Npn \__zhvt_set_shukou_content: {
	% ★ 鱼尾上方留白（title 静态框）嵌入闲叟钤印：
	%   整页 /Rotate 90，故此处把图用 \rotatebox{90} 转正为竖章；
	%   尺寸卡在版心条幅内（书内约 7mm 宽 × 33.8mm 高，留余量不压鱼尾）。
	\setstaticcontents * { title } {
		\vbox_to_ht:nn { \l__zhvt_baseline_skip_dim } {
			\vskip 14.2pt
			\hbox_to_wd:nn { \l__zhvt_shukou_title_height_dim } {
				\hfil
				\rotatebox [ origin = c ] { 90 } { \includegraphics [ width = 7mm , height = 33.8mm ] { yshs.png } }
				\hfil
			}
		}
	}
% ★ 版心（书口）统一改法（2026-11-16 愚大人定）：18pt＋框内居中＋书口文字同版框色
%   ① 原 27.2pt 时 7 字长篇名（如「廉颇蔺相如列传」）末字溢出 chapter 框、压住页码；
%      改 \fontsize{18pt}{18pt}，前后 \hfil 令字串在框内居中（不随篇名字数漂移）。
%   ② \vbox{\vskip 8.2pt} 定量下移，抵消缩字号后栏宽方向 8.2pt 上偏（像素实测 y 中心 543.3 = 原版）。
%      ⚠ \raisebox / \vspace* 在 flowfram 框内无效；换字号须同步换补偿：16pt→10.1、19pt→7.4。
%   ③ chapter 与 page~number 均置 \color{sealbrown}（HTML 874434），与版框／界行同色；正文仍黑。
%   ⚠ 页码字位、钤印、鱼尾、界行、正文一律未动；五部合订本同此口径。
	\setdynamiccontents * { chapter }
	{ \color{sealbrown} \vbox { \vskip 8.2pt \hbox to \l__zhvt_shukou_chapter_title_height_dim
		{ \hfil \fontsize{18pt}{18pt}\selectfont \l__zhvt_title_str \quad \l__zhvt_chapter_title_tl \hfil } } }
	\setdynamiccontents * { page~number } { \color{sealbrown} \quad \zhnum { page } }
}
\__zhvt_set_shukou_content:
\ExplSyntaxOff

\maketitle

% 单篇，无多卷目录，不渲染目录页

\juan{淮南衡山列传}{卷一百十八~淮南衡山列传第五十八}
\hspace*{2em}淮南厉王长者。高祖少子也。其母故赵王张敖美人。高祖八年。从东垣过赵。赵王献之美人。厉王母得幸焉。有身。赵王敖弗敢内宫。为筑外宫而舍之。及贯高等谋反柏人事发觉。并逮治王。尽收捕王母兄弟美人。系之河内。厉王母亦系。告吏曰。﹁得幸上。有身。﹂吏以闻上。上方怒赵王。未理厉王母。厉王母弟赵兼因辟阳侯言吕后。吕后妒。弗肯白。辟阳侯不彊争。及厉王母已生厉王。恚。即自杀。吏奉厉王诣上。上悔。令吕后母之。而葬厉王母真定。真定。厉王母之家在焉。父世县也。\\
\hspace*{2em}高祖十一年月。淮南王黥布反。立子长为淮南王。王黥布故地。凡四郡。上自将兵击灭布。厉王遂即位。厉王蚤失母。常附吕后。孝惠。吕后时以故得幸无患害。而常心怨辟阳侯。弗敢发。及孝文帝初即位。淮南王自以为最亲。骄蹇。数不奉法。上以亲故。常宽赦之。三年。入朝。甚横。从上入苑囿猎。与上同车。常谓上﹁大兄﹂。厉王有材力。力能扛鼎。乃往请辟阳侯。辟阳侯出见之。即自袖铁椎椎辟阳侯。令从者魏敬刭之。厉王乃驰走阙下。肉袒谢曰。﹁臣母不当坐赵事。其时辟阳侯力能得之吕后。弗争。罪一也。赵王如意子母无罪。吕后杀之。辟阳侯弗争。罪二也。吕后王诸吕。欲以危刘氏。辟阳侯弗争。罪三也。臣谨为天下诛贼臣辟阳侯。报母之仇。谨伏阙下请罪。﹂孝文伤其志。为亲故。弗治。赦厉王。当是时。薄太后及太子诸大臣皆惮厉王。厉王以此归国益骄恣。不用汉法。出入称警跸。称制。自为法令。拟於天子。\\
\hspace*{2em}六年。令男子但等七十人与棘蒲侯柴武太子奇谋。以輂车四十乘反谷口。令人使闽越。匈奴。事觉。治之。使使召淮南王。淮南王至长安。\\
\hspace*{2em}﹁丞相臣张仓。典客臣冯敬。行御史大夫事宗正臣逸。廷尉臣贺。备盗贼中尉臣福昧死言。淮南王长废先帝法。不听天子诏。居处无度。为黄屋盖乘舆。出入拟於天子。擅为法令。不用汉法。及所置吏。以其郎中春为丞相。聚收汉诸侯人及有罪亡者。匿与居。为治家室。赐其财物爵禄田宅。爵或至关内侯。奉以二千石。所不当得。欲以有为。大夫但。士五开章等七十人与棘蒲侯太子奇谋反。欲以危宗庙社稷。使开章阴告长。与谋使闽越及匈奴发其兵。开章之淮南见长。长数与坐语饮食。为家室娶妇。以二千石俸奉之。开章使人告但。已言之王。春使使报但等。吏觉知。使长安尉奇等往捕开章。长匿不予。与故中尉蕑忌谋。杀以闭口。为棺椁衣衾。葬之肥陵邑。谩吏曰﹃不知安在﹄。又详聚土。树表其上。曰﹃开章死。埋此下﹄。及长身自贼杀无罪者一人。令吏论杀无罪者六人。为亡命弃市罪诈捕命者以除罪。擅罪人。罪人无告劾。系治城旦舂以上十四人。赦免罪人。死罪十八人。城旦舂以下五十八人。赐人爵关内侯以下九十四人。前日长病。陛下忧苦之。使使者赐书。枣脯。长不欲受赐。不肯见拜使者。南海民处庐江界中者反。淮南吏卒击之。陛下以淮南民贫苦。遣使者赐长帛五千匹。以赐吏卒劳苦者。长不欲受赐。谩言曰﹃无劳苦者﹄。南海民王织上书献璧皇帝。忌擅燔其书。不以闻。吏请召治忌。长不遣。谩言曰﹃忌病﹄。春又请长。原入见。长怒曰﹃女欲离我自附汉﹄。长当弃市。臣请论如法。﹂\\
\hspace*{2em}制曰。﹁朕不忍致法於王。其与列侯二千石议。﹂\\
\hspace*{2em}﹁臣仓。臣敬。臣逸。臣福。臣贺昧死言。臣谨与列侯吏二千石臣婴等四十三人议。皆曰﹃长不奉法度。不听天子诏。乃阴聚徒党及谋反者。厚养亡命。欲以有为﹄。臣等议论如法。﹂\\
\hspace*{2em}制曰。﹁朕不忍致法於王。其赦长死罪。废勿王。﹂\\
\hspace*{2em}﹁臣仓等昧死言。长有大死罪。陛下不忍致法。幸赦。废勿王。臣请处蜀郡严道邛邮。遣其子母从居。县为筑盖家室。皆廪食给薪菜盐豉炊食器席蓐。臣等昧死请。请布告天下。﹂\\
\hspace*{2em}制曰。﹁计食长给肉日五斤。酒二斗。令故美人才人得幸者十人从居。他可。﹂\\
\hspace*{2em}尽诛所与谋者。於是乃遣淮南王。载以辎车。令县以次传。是时袁盎谏上曰。﹁上素骄淮南王。弗为置严傅相。以故至此。且淮南王为人刚。今暴摧折之。臣恐卒逢雾露病死。陛下为有杀弟之名。柰何。﹂上曰。﹁吾特苦之耳。今复之。﹂县传淮南王者皆不敢发车封。淮南王乃谓侍者曰。﹁谁谓乃公勇者。吾安能勇。吾以骄故不闻吾过至此。人生一世间。安能邑邑如此。﹂乃不食死。至雍。雍令发封。以死闻。上哭甚悲。谓袁盎曰。﹁吾不听公言。卒亡淮南王。﹂盎曰。﹁不可柰何。原陛下自宽。﹂上曰。﹁为之柰何。﹂盎曰。﹁独斩丞相。御史以谢天下乃可。﹂上即令丞相。御史逮考诸县传送淮南王不发封餽侍者。皆弃市。乃以列侯葬淮南王於雍。守冢三十户。\\
\hspace*{2em}孝文八年。上怜淮南王。淮南王有子四人。皆七八岁。乃封子安为阜陵侯。子勃为安阳侯。子赐为阳周侯。子良为东成侯。\\
\hspace*{2em}孝文十二年。民有作歌歌淮南厉王曰。﹁一尺布。尚可缝。一斗粟。尚可舂。兄弟二人不能相容。﹂上闻之。乃叹曰。﹁尧舜放逐骨肉。周公杀管蔡。天下称圣。何者。不以私害公。天下岂以我为贪淮南王地邪。﹂乃徙城阳王王淮南故地。而追尊谥淮南王为厉王。置园复如诸侯仪。\\
\hspace*{2em}孝文十六年。徙淮南王喜复故城阳。上怜淮南厉王废法不轨。自使失国蚤死。乃立其三子。阜陵侯安为淮南王。安阳侯勃为衡山王。阳周侯赐为庐江王。皆复得厉王时地。参分之。东城侯良前薨。无后也。\\
\hspace*{2em}孝景三年。吴楚七国反。吴使者至淮南。淮南王欲发兵应之。其相曰。﹁大王必欲发兵应吴。臣原为将。﹂王乃属相兵。淮南相已将兵。因城守。不听王而为汉。汉亦使曲城侯将兵救淮南。淮南以故得完。吴使者至庐江。庐江王弗应。而往来使越。吴使者至衡山。衡山王坚守无二心。孝景四年。吴楚已破。衡山王朝。上以为贞信。乃劳苦之曰。﹁南方卑湿。﹂徙衡山王王济北。所以襃之。及薨。遂赐谥为贞王。庐江王边越。数使使相交。故徙为衡山王。王江北。淮南王如故。\\
\hspace*{2em}淮南王安为人好读书鼓琴。不喜弋猎狗马驰骋。亦欲以行阴德拊循百姓。流誉天下。时时怨望厉王死。时欲畔逆。未有因也。及建元二年。淮南王入朝。素善武安侯。武安侯时为太尉。乃逆王霸上。与王语曰。﹁方今上无太子。大王亲高皇帝孙。行仁义。天下莫不闻。即宫车一日晏驾。非大王当谁立者。﹂淮南王大喜。厚遗武安侯金财物。阴结宾客。拊循百姓。为畔逆事。建元六年。彗星见。淮南王心怪之。或说王曰。﹁先吴军起时。彗星出长数尺。然尚流血千里。今彗星长竟天。天下兵当大起。﹂王心以为上无太子。天下有变。诸侯并争。愈益治器械攻战具。积金钱赂遗郡国诸侯游士奇材。诸辨士为方略者。妄作妖言。谄谀王。王喜。多赐金钱。而谋反滋甚。\\
\hspace*{2em}淮南王有女陵。慧。有口辩。王爱陵。常多予金钱。为中诇长安。约结上左右。元朔三年。上赐淮南王几杖。不朝。淮南王王后荼。王爱幸之。王后生太子迁。迁取王皇太后外孙修成君女为妃。王谋为反具。畏太子妃知而内泄事。乃与太子谋。令诈弗爱。三月不同席。王乃详为怒太子。闭太子使与妃同内三月。太子终不近妃。妃求去。王乃上书谢归去之。王后荼。太子迁及女陵得爱幸王。擅国权。侵夺民田宅。妄致系人。\\
\hspace*{2em}元朔五年。太子学用剑。自以为人莫及。闻郎中雷被巧。乃召与戏。被一再辞让。误中太子。太子怒。被恐。此时有欲从军者辄诣京师。被即原奋击匈奴。太子迁数恶被於王。王使郎中令斥免。欲以禁后。被遂亡至长安。上书自明。诏下其事廷尉。河南。河南治。逮淮南太子。王。王后计欲无遣太子。遂发兵反。计犹豫。十馀日未定。会有诏。即讯太子。当是时。淮南相怒寿春丞留太子逮不遣。劾不敬。王以请相。相弗听。王使人上书告相。事下廷尉治。踪迹连王。王使人候伺汉公卿。公卿请逮捕治王。王恐事发。太子迁谋曰。﹁汉使即逮王。王令人衣卫士衣。持戟居庭中。王旁有非是。则刺杀之。臣亦使人刺杀淮南中尉。乃举兵。未晚。﹂是时上不许公卿请。而遣汉中尉宏即讯验王。王闻汉使来。即如太子谋计。汉中尉至。王视其颜色和。讯王以斥雷被事耳。王自度无何。不发。中尉还。以闻。公卿治者曰。﹁淮南王安拥阏奋击匈奴者雷被等。废格明诏。当弃市。﹂诏弗许。公卿请废勿王。诏弗许。公卿请削五县。诏削二县。使中尉宏赦淮南王罪。罚以削地。中尉入淮南界。宣言赦王。王初闻汉公卿请诛之。未知得削地。闻汉使来。恐其捕之。乃与太子谋刺之如前计。及中尉至。即贺王。王以故不发。其后自伤曰。﹁吾行仁义见削。甚耻之。﹂然淮南王削地之后。其为反谋益甚。诸使道从长安来。为妄妖言。言上无男。汉不治。即喜。即言汉廷治。有男。王怒。以为妄言。非也。\\
\hspace*{2em}王日夜与伍被。左吴等案舆地图。部署兵所从入。王曰。﹁上无太子。宫车即晏驾。廷臣必徵胶东王。不即常山王。诸侯并争。吾可以无备乎。且吾高祖孙。亲行仁义。陛下遇我厚。吾能忍之。万世之后。吾宁能北面臣事竖子乎。﹂\\
\hspace*{2em}王坐东宫。召伍被与谋。曰。﹁将军上。﹂被怅然曰。﹁上宽赦大王。王复安得此亡国之语乎。臣闻子胥谏吴王。吴王不用。乃曰﹃臣今见麋鹿游姑苏之台也﹄。今臣亦见宫中生荆棘。露霑衣也。﹂王怒。系伍被父母。囚之三月。复召曰。﹁将军许寡人乎。﹂被曰。﹁不。直来为大王画耳。臣闻聪者听於无声。明者见於未形。故圣人万举万全。昔文王一动而功显于千世。列为三代。此所谓因天心以动作者也。故海内不期而随。此千岁之可见者。夫百年之秦。近世之吴楚。亦足以喻国家之存亡矣。臣不敢避子胥之诛。原大王毋为吴王之听。昔秦绝圣人之道。杀术士。燔诗书。弃礼义。尚诈力。任刑罚。转负海之粟致之西河。当是之时。男子疾耕不足於糟糠。女子纺绩不足於盖形。遣蒙恬筑长城。东西数千里。暴兵露师常数十万。死者不可胜数。僵尸千里。流血顷亩。百姓力竭。欲为乱者十家而五。又使徐福入海求神异物。还为伪辞曰。﹃臣见海中大神。言曰。﹁汝西皇之使邪。﹂臣答曰。﹁然。﹂﹁汝何求。﹂曰。﹁原请延年益寿药。﹂神曰。﹁汝秦王之礼薄。得观而不得取。﹂即从臣东南至蓬莱山。见芝成宫阙。有使者铜色而龙形。光上照天。於是臣再拜问曰。﹁宜何资以献。﹂海神曰。﹁以令名男子若振女与百工之事。即得之矣。﹂﹄秦皇帝大说。遣振男女三千人。资之五穀种种百工而行。徐福得平原广泽。止王不来。於是百姓悲痛相思。欲为乱者十家而六。又使尉佗逾五岭攻百越。尉佗知中国劳极。止王不来。使人上书。求女无夫家者三万人。以为士卒衣补。秦皇帝可其万五千人。於是百姓离心瓦解。欲为乱者十家而七。客谓高皇帝曰。﹃时可矣。﹄高皇帝曰。﹃待之。圣人当起东南间。﹄不一年。陈胜吴广发矣。高皇始於丰沛。一倡天下不期而响应者不可胜数也。此所谓蹈瑕候间。因秦之亡而动者也。百姓原之。若旱之望雨。故起於行陈之中而立为天子。功高三王。德传无穷。今大王见高皇帝得天下之易也。独不观近世之吴楚乎。夫吴王赐号为刘氏祭酒。复不朝。王四郡之众。地方数千里。内铸消铜以为钱。东煮海水以为盐。上取江陵木以为船。一船之载当中国数十两车。国富民众。行珠玉金帛赂诸侯宗室大臣。独窦氏不与。计定谋成。举兵而西。破於大梁。败於狐父。奔走而东。至於丹徒。越人禽之。身死绝祀。为天下笑。夫以吴越之众不能成功者何。诚逆天道而不知时也。方今大王之兵众不能十分吴楚之一。天下安宁有万倍於秦之时。原大王从臣之计。大王不从臣之计。今见大王事必不成而语先泄也。臣闻微子过故国而悲。於是作麦秀之歌。是痛纣之不用王子比干也。故孟子曰﹃纣贵为天子。死曾不若匹夫﹄。是纣先自绝於天下久矣。非死之日而天下去之。今臣亦窃悲大王弃千乘之君。必且赐绝命之书。为群臣先。死於东宫也。﹂於是气怨结而不扬。涕满匡而横流。即起。历阶而去。\\
\hspace*{2em}王有孽子不害。最长。王弗爱。王。王后。太子皆不以为子兄数。不害有子建。材高有气。常怨望太子不省其父。又怨时诸侯皆得分子弟为侯。而淮南独二子。一为太子。建父独不得为侯。建阴结交。欲告败太子。以其父代之。太子知之。数捕系而榜笞建。建具知太子之谋欲杀汉中尉。即使所善寿春庄芷以元朔六年上书於天子曰。﹁毒药苦於口利於病。忠言逆於耳利於行。今淮南王孙建。材能高。淮南王王后荼。荼子太子迁常疾害建。建父不害无罪。擅数捕系。欲杀之。今建在。可徵问。具知淮南阴事。﹂书闻。上以其事下廷尉。廷尉下河南治。是时故辟阳侯孙审卿善丞相公孙弘。怨淮南厉王杀其大父。乃深购淮南事於弘。弘乃疑淮南有畔逆计谋。深穷治其狱。河南治建。辞引淮南太子及党与。淮南王患之。欲发。问伍被曰。﹁汉廷治乱。﹂伍被曰。﹁天下治。﹂王意不说。谓伍被曰。﹁公何以言天下治也。﹂被曰。﹁被窃观朝廷之政。君臣之义。父子之亲。夫妇之别。长幼之序。皆得其理。上之举错遵古之道。风俗纪纲未有所缺也。重装富贾。周流天下。道无不通。故交易之道行。南越宾服。羌僰入献。东瓯入降。广长榆。开朔方。匈奴折翅伤翼。失援不振。虽未及古太平之时。然犹为治也。﹂王怒。被谢死罪。王又谓被曰。﹁山东即有兵。汉必使大将军将而制山东。公以为大将军何如人也。﹂被曰。﹁被所善者黄义。从大将军击匈奴。还。告被曰。﹃大将军遇士大夫有礼。於士卒有恩。众皆乐为之用。骑上下山若蜚。材幹绝人。﹄被以为材能如此。数将习兵。未易当也。及谒者曹梁使长安来。言大将军号令明。当敌勇敢。常为士卒先。休舍。穿井未通。须士卒尽得水。乃敢饮。军罢。卒尽已度河。乃度。皇太后所赐金帛。尽以赐军吏。虽古名将弗过也。﹂王默然。\\
\hspace*{2em}淮南王见建已徵治。恐国阴事且觉。欲发。被又以为难。乃复问被曰。﹁公以为吴兴兵是邪非也。﹂被曰。﹁以为非也。吴王至富贵也。举事不当。身死丹徒。头足异处。子孙无遗类。臣闻吴王悔之甚。原王孰虑之。无为吴王之所悔。﹂王曰。﹁男子之所死者一言耳。且吴何知反。汉将一日过成皋者四十馀人。今我令楼缓先要成皋之口。周被下颍川兵塞轘辕。伊阙之道。陈定发南阳兵守武关。河南太守独有雒阳耳。何足忧。然此北尚有临晋关。河东。上党与河内。赵国。人言曰﹃绝成皋之口。天下不通﹄。据三川之险。招山东之兵。举事如此。公以为何如。﹂被曰。﹁臣见其祸。未见其福也。﹂王曰。﹁左吴。赵贤。硃骄如皆以为有福。什事九成。公独以为有祸无福。何也。﹂被曰。﹁大王之群臣近幸素能使众者。皆前系诏狱。馀无可用者。﹂王曰。﹁陈胜。吴广无立锥之地。千人之聚。起於大泽。奋臂大呼而天下响应。西至於戏而兵百二十万。今吾国虽小。然而胜兵者可得十馀万。非直適戍之众。釠凿棘矜也。公何以言有祸无福。﹂被曰。﹁往者秦为无道。残贼天下。兴万乘之驾。作阿房之宫。收太半之赋。发闾左之戍。父不宁子。兄不便弟。政苛刑峻。天下熬然若焦。民皆引领而望。倾耳而听。悲号仰天。叩心而怨上。故陈胜大呼。天下响应。当今陛下临制天下。一齐海内。汎爱蒸庶。布德施惠。口虽未言。声疾雷霆。令虽未出。化驰如神。心有所怀。威动万里。下之应上。犹影响也。而大将军材能不特章邯。杨熊也。大王以陈胜。吴广谕之。被以为过矣。﹂王曰。﹁苟如公言。不可徼幸邪。﹂被曰。﹁被有愚计。﹂王曰。﹁柰何。﹂被曰。﹁当今诸侯无异心。百姓无怨气。朔方之郡田地广。水草美。民徙者不足以实其地。臣之愚计。可伪为丞相御史请书。徙郡国豪桀任侠及有耐罪以上。赦令除其罪。产五十万以上者。皆徙其家属朔方之郡。益发甲卒。急其会日。又伪为左右都司空上林中都官诏狱书。诸侯太子幸臣。如此则民怨。诸侯惧。即使辩武随而说之。傥可徼幸什得一乎。﹂王曰。﹁此可也。虽然。吾以为不至若此。﹂於是王乃令官奴入宫。作皇帝玺。丞相。御史。大将军。军吏。中二千石。都官令。丞印。及旁近郡太守。都尉印。汉使节法冠。欲如伍被计。使人伪得罪而西。事大将军。丞相。一日发兵。使人即刺杀大将军青。而说丞相下之。如发蒙耳。\\
\hspace*{2em}王欲发国中兵。恐其相。二千石不听。王乃与伍被谋。先杀相。二千石。伪失火宫中。相。二千石救火。至即杀之。计未决。又欲令人衣求盗衣。持羽檄。从东方来。呼曰﹁南越兵入界﹂。欲因以发兵。乃使人至庐江。会稽为求盗。未发。王问伍被曰。﹁吾举兵西乡。诸侯必有应我者。即无应。柰何。﹂被曰。﹁南收衡山以击庐江。有寻阳之船。守下雉之城。结九江之浦。绝豫章之口。彊弩临江而守。以禁南郡之下。东收江都。会稽。南通劲越。屈彊江淮间。犹可得延岁月之寿。﹂王曰。﹁善。无以易此。急则走越耳。﹂\\
\hspace*{2em}於是廷尉以王孙建辞连淮南王太子迁闻。上遣廷尉监因拜淮南中尉。逮捕太子。至淮南。淮南王闻。与太子谋召相。二千石。欲杀而发兵。召相。相至。内史以出为解。中尉曰。﹁臣受诏使。不得见王。﹂王念独杀相而内史中尉不来。无益也。即罢相。王犹豫。计未决。太子念所坐者谋刺汉中尉。所与谋者已死。以为口绝。乃谓王曰。﹁群臣可用者皆前系。今无足与举事者。王以非时发。恐无功。臣原会逮。﹂王亦偷欲休。即许太子。太子即自刭。不殊。伍被自诣吏。因告与淮南王谋反。反踪迹具如此。\\
\hspace*{2em}吏因捕太子。王后。围王宫。尽求捕王所与谋反宾客在国中者。索得反具以闻。上下公卿治。所连引与淮南王谋反列侯二千石豪杰数千人。皆以罪轻重受诛。衡山王赐。淮南王弟也。当坐收。有司请逮捕衡山王。天子曰。﹁诸侯各以其国为本。不当相坐。与诸侯王列侯会肄丞相诸侯议。﹂赵王彭祖。列侯臣让等四十三人议。皆曰。﹁淮南王安甚大逆无道。谋反明白。当伏诛。﹂胶西王臣端议曰。﹁淮南王安废法行邪。怀诈伪心。以乱天下。荧惑百姓。倍畔宗庙。妄作妖言。春秋曰﹃臣无将。将而诛﹄。安罪重於将。谋反形已定。臣端所见其书节印图及他逆无道事验明白。甚大逆无道。当伏其法。而论国吏二百石以上及比者。宗室近幸臣不在法中者。不能相教。当皆免官削爵为士伍。毋得宦为吏。其非吏。他赎死金二斤八两。以章臣安之罪。使天下明知臣子之道。毋敢复有邪僻倍畔之意。﹂丞相弘。廷尉汤等以闻。天子使宗正以符节治王。未至。淮南王安自刭杀。王后荼。太子迁诸所与谋反者皆族。天子以伍被雅辞多引汉之美。欲勿诛。廷尉汤曰。﹁被首为王画反谋。被罪无赦。﹂遂诛被。国除为九江郡。\\
\hspace*{2em}衡山王赐。王后乘舒生子三人。长男爽为太子。次男孝。次女无采。又姬徐来生子男女四人。美人厥姬生子二人。衡山王。淮南王兄弟相责望礼节。间不相能。衡山王闻淮南王作为畔逆反具。亦心结宾客以应之。恐为所并。\\
\hspace*{2em}元光六年。衡山王入朝。其谒者卫庆有方术。欲上书事天子。王怒。故劾庆死罪。彊榜服之。衡山内史以为非是。卻其狱。王使人上书告内史。内史治。言王不直。王又数侵夺人田。坏人冢以为田。有司请逮治衡山王。天子不许。为置吏二百石以上。衡山王以此恚。与奚慈。张广昌谋。求能为兵法候星气者。日夜从容王密谋反事。\\
\hspace*{2em}王后乘舒死。立徐来为王后。厥姬俱幸。两人相妒。厥姬乃恶王后徐来於太子曰。﹁徐来使婢蛊道杀太子母。﹂太子心怨徐来。徐来兄至衡山。太子与饮。以刃刺伤王后兄。王后怨怒。数毁恶太子於王。太子女弟无采。嫁弃归。与奴奸。又与客奸。太子数让无采。无采怒。不与太子通。王后闻之。即善遇无采。无采及中兄孝少失母。附王后。王后以计爱之。与共毁太子。王以故数击笞太子。元朔四年中。人有贼伤王后假母者。王疑太子使人伤之。笞太子。后王病。太子时称病不侍。孝。王后。无采恶太子。﹁太子实不病。自言病。有喜色。﹂王大怒。欲废太子。立其弟孝。王后知王决废太子。又欲并废孝。王后有侍者。善舞。王幸之。王后欲令侍者与孝乱以汙之。欲并废兄弟而立其子广代太子。太子爽知之。念后数恶己无已时。欲与乱以止其口。王后饮。太子前为寿。因据王后股。求与王后卧。王后怒。以告王。王乃召。欲缚而笞之。太子知王常欲废己立其弟孝。乃谓王曰。﹁孝与王御者奸。无采与奴奸。王彊食。请上书。﹂即倍王去。王使人止之。莫能禁。乃自驾追捕太子。太子妄恶言。王械系太子宫中。孝日益亲幸。王奇孝材能。乃佩之王印。号曰将军。令居外宅。多给金钱。招致宾客。宾客来者。微知淮南。衡山有逆计。日夜从容劝之。王乃使孝客江都人救赫。陈喜作輣车镞矢。刻天子玺。将相军吏印。王日夜求壮士如周丘等。数称引吴楚反时计画。以约束。衡山王非敢效淮南王求即天子位。畏淮南起并其国。以为淮南已西。发兵定江淮之间而有之。望如是。\\
\hspace*{2em}元朔五年秋。衡山王当朝。过淮南。淮南王乃昆弟语。除前卻。约束反具。衡山王即上书谢病。上赐书不朝。\\
\hspace*{2em}元朔六年中。衡山王使人上书请废太子爽。立孝为太子。爽闻。即使所善白嬴之长安上书。言孝作輣车镞矢。与王御者奸。欲以败孝。白嬴至长安。未及上书。吏捕嬴。以淮南事系。王闻爽使白嬴上书。恐言国阴事。即上书反告太子爽所为不道弃市罪事。事下沛郡治。元年冬。有司公卿下沛郡求捕所与淮南谋反者未得。得陈喜於衡山王子孝家。吏劾孝首匿喜。孝以为陈喜雅数与王计谋反。恐其发之。闻律先自告除其罪。又疑太子使白嬴上书发其事。即先自告。告所与谋反者救赫。陈喜等。廷尉治验。公卿请逮捕衡山王治之。天子曰。﹁勿捕。﹂遣中尉安。大行息即问王。王具以情实对。吏皆围王宫而守之。中尉大行还。以闻。公卿请遣宗正。大行与沛郡杂治王。王闻。即自刭杀。孝先自告反。除其罪。坐与王御婢奸。弃市。王后徐来亦坐蛊杀前王后乘舒。及太子爽坐王告不孝。皆弃市。诸与衡山王谋反者皆族。国除为衡山郡。\\
\hspace*{2em}太史公曰。诗之所谓﹁戎狄是膺。荆舒是惩﹂。信哉是言也。淮南。衡山亲为骨肉。疆土千里。列为诸侯。不务遵蕃臣职以承辅天子。而专挟邪僻之计。谋为畔逆。仍父子再亡国。各不终其身。为天下笑。此非独王过也。亦其俗薄。臣下渐靡使然也。夫荆楚僄勇轻悍。好作乱。乃自古记之矣。\\

% ===== 索隐述赞（司马贞《史记索隐》）· 统一样式 =====
% 题头「索隐述赞曰」五字并入反白框；正文每四句一列，定宽盒 \hfil 居列中；
% \fontsize{24}{41} 列距对齐全书界行（121+41k）；\clearpage 使述赞独占一页。
\clearpage
\shuzanhead{索隐述赞曰}\linebreak
{\shuzanbody\fontsize{24}{41}\selectfont
\hbox to \hsize{\hfil 淮南多横。举事非正。天子宽仁。其过不更。\hfil}\linebreak
\hbox to \hsize{\hfil 轞车致祸。斗粟成咏。王安好学。女陵作诇。\hfil}\linebreak
\hbox to \hsize{\hfil 兄弟不和。倾国殒命。\hfil}%
\par}
\end{document}
